\documentclass[11pt,a4paper]{article}

% ------------------------------------------------
% Pacotes básicos
% ------------------------------------------------

\usepackage[utf8]{inputenc}
\usepackage[T1]{fontenc}
\usepackage[brazil]{babel}

\usepackage{amsmath}
\usepackage{amssymb}
\usepackage{graphicx}
\usepackage{booktabs}
\usepackage{geometry}

\geometry{
    margin=2.5cm
}


% ------------------------------------------------
% Informações do documento
% ------------------------------------------------

\title{
Geometria Espectral Relacional\\
Série S26-B\\
Regime Não Linear e Auditorias Dinâmicas
}

\author{
Eduardo Batista de Freitas 
}

\date{\today}


% ------------------------------------------------
% Documento
% ------------------------------------------------

\begin{document}

\maketitle


\begin{abstract}

Este documento apresenta a continuidade experimental da
Geometria Espectral Relacional após a validação numérica
do regime linear.

A série S26-B investiga a introdução controlada de termos
não lineares, a estabilidade dinâmica, a redistribuição
espectral de energia e a possível existência de estados
estacionários do operador relacional.

\end{abstract}



\tableofcontents


\section{Introdução}

A série S26-B representa a transição entre a validação
do regime linear e a investigação sistemática da dinâmica
não linear.

O ciclo anterior estabeleceu a confiabilidade do motor
computacional através da conservação Hamiltoniana,
convergência temporal e validação espectral.

A presente etapa mantém a mesma estrutura relacional,
introduzindo parâmetros não lineares controlados.


\section{Objetivos}

Os objetivos principais são:

\begin{itemize}

\item determinar regimes estáveis de evolução não linear;

\item caracterizar redistribuição de energia no espaço espectral;

\item medir taxas de mistura modal;

\item investigar estados persistentes;

\item analisar dependência topológica futura.

\end{itemize}



\section{Metodologia}

A evolução permanece baseada no operador relacional
discreto e sua decomposição espectral.

A dinâmica é integrada utilizando o mesmo formalismo
numérico validado na etapa linear.


\section{Resultados}

Esta seção será preenchida após a consolidação das
auditorias S26-B.

\section{Auditoria 12 — Robustez Temporal em Regime Não Linear}

A robustez temporal do integrador foi avaliada no regime
fortemente não linear, utilizando $\beta=100$.

Foram comparados diferentes passos temporais:

\[
\Delta t =
5\times10^{-5},
2.5\times10^{-5},
1.25\times10^{-5},
6.25\times10^{-6}
\]

Observou-se redução sistemática do erro relativo de energia
com o refinamento temporal, indicando convergência compatível
com um integrador de segunda ordem.

\begin{table}[h]
\centering
\begin{tabular}{ccccc}
\toprule
$\Delta t$ & Erro & Amplitude & Entropia & Largura\\
\midrule
$5e-5$ & 3.58 & 1.678 & 2.913 & 11.011\\
$2.5e-5$ & 0.384 & 1.133 & 2.765 & 9.136\\
$1.25e-5$ & 0.080 & 1.031 & 2.731 & 8.767\\
$6.25e-6$ & 0.019 & 1.007 & 2.723 & 8.681\\
\bottomrule
\end{tabular}
\end{table}

O resultado indica que as alterações observadas no regime
fortemente não linear para passos temporais maiores são
efeitos de discretização e não mudanças estruturais do sistema.

\section{Auditoria 13 — Dependência com o parâmetro $\beta$}

Foi realizada uma análise paramétrica do comportamento
dinâmico em função do parâmetro não linear $\beta$.

Foram avaliados os valores:

\[
\beta =
0.1,\ 0.25,\ 0.5,\ 1.0,\ 2.0
\]

Os resultados obtidos foram:

\begin{table}[h]
\centering
\begin{tabular}{ccccc}
\toprule
$\beta$ & Energia & Amplitude & Modo Dominante & PR\\
\midrule
0.1 & -2.55e-2 & 1.0126 & 0 & 1.00003\\
0.25 & -6.76e-2 & 1.0318 & 0 & 1.00003\\
0.5 & -1.49e-1 & 1.0646 & 0 & 1.00003\\
1.0 & -3.65e-1 & 1.1336 & 0 & 1.00004\\
2.0 & -1.13 & 1.2871 & 0 & 1.00006\\
\bottomrule
\end{tabular}
\end{table}


Observou-se uma redução monotônica da energia com o
aumento de $\beta$:

\[
\frac{\partial E}{\partial \beta}<0
\]

no intervalo analisado.

A amplitude final apresentou crescimento progressivo,
enquanto o modo espectral dominante permaneceu invariável:

\[
m_{\mathrm{dominante}}=0
\]

O índice de participação modal permaneceu próximo de:

\[
PR \approx 1
\]

indicando ausência de mistura modal significativa no regime
avaliado.

\section{Discussão}

A interpretação dos resultados será realizada somente
após a conclusão das auditorias previstas.


\section{Conclusão}

A série S26-B encontra-se em desenvolvimento.

% Remova a linha antiga e adicione uma quebra de página 
% para o Journal começar perfeitamente como uma nova parte do documento
\section{Auditoria 13.1 — Análise de Sensibilidade ao Parâmetro $\beta$}

A análise da dependência com o parâmetro não linear
$\beta$ foi realizada utilizando os resultados da Auditoria 13.

No intervalo:

\[
0.1 \leq \beta \leq 2.0
\]

observou-se redução monotônica da energia:

\[
\frac{\partial E}{\partial\beta}<0
\]

para todos os pontos avaliados.

A amplitude final apresentou aumento relativo de:

\[
\frac{\Delta A}{A_0}=27.106\%
\]

mantendo, entretanto, a estrutura espectral estável.

O modo dominante permaneceu invariável:

\[
m_{\mathrm{dominante}}=0
\]

em toda a faixa analisada.

O índice de participação modal apresentou pequena variação:

\[
\Delta PR =
3.29\times10^{-5}
\]

com valor médio:

\[
\overline{PR}=1.000039
\]

indicando ausência de redistribuição modal significativa.
\section{Auditoria 14 — Transição de Regime Não Linear}

A varredura inicial do parâmetro $\beta$ indicou
estabilidade espectral até aproximadamente
$\beta=10$.

Nesse intervalo:

\[
m_{\mathrm{dominante}}=0
\]

permaneceu constante e o índice de participação modal
apresentou pequena variação.

Para valores superiores observou-se crescimento extremo
da energia e ocorrência de overflow numérico, indicando
necessidade de refinamento temporal antes de qualquer
interpretação estrutural.
\section{Auditoria 13 — Dependência entre Escala Não Linear e Resolução Temporal}

A análise de estabilidade S26-B24 investigou a relação entre
a intensidade não linear do sistema, representada pelo parâmetro
$\beta$, e o passo temporal de integração $\Delta t$.

Os resultados mostram que regimes aparentemente divergentes
para passos temporais elevados recuperam estabilidade quando
a resolução temporal é refinada.

Observou-se uma dependência qualitativa:

\[
\beta \uparrow
\Rightarrow
\Delta t_{crit} \downarrow
\]

indicando a existência de uma escala temporal crítica associada
à intensidade dinâmica do operador relacional.

O fenômeno sugere que a variável temporal não atua apenas como
parâmetro externo de evolução, mas pode estar acoplada à própria
estrutura dinâmica do sistema.

Para $\beta=20$, por exemplo, o regime
$\Delta t=2.5\times10^{-4}$ apresentou divergência,
enquanto os passos

\[
1.25\times10^{-4},
6.25\times10^{-5},
3.125\times10^{-5}
\]

mantiveram estabilidade.

Esse comportamento indica que parte das instabilidades observadas
em regimes fortemente não lineares são associadas à resolução
temporal insuficiente e não necessariamente à ausência de estados
dinâmicos estáveis.
\begin{table}[h]
\centering
\begin{tabular}{cccc}
\toprule
$\beta$ & $\Delta t$ crítico aproximado & Estável & Observação\\
\midrule
10 & 2.5e-4 & Sim & Alta deformação\\
20 & 1.25e-4 & Sim & Transição\\
30 & 1.25e-4 & Sim & Forte não linearidade\\
50 & 1.25e-4 & Sim & Regime extremo\\
\bottomrule
\end{tabular}
\end{table}
\section{Auditoria 14 — Escala Crítica Temporal}

A análise da curva crítica temporal investigou a dependência
entre o parâmetro não linear $\beta$ e o maior passo temporal
compatível com estabilidade dinâmica.

Os resultados indicam uma redução sistemática do limite temporal
com o aumento da intensidade não linear:

\[
\beta \uparrow
\Rightarrow
\Delta t_{crit}\downarrow
\]

A fronteira observada apresenta comportamento em patamares:

\[
5\times10^{-4}
\rightarrow
2.5\times10^{-4}
\rightarrow
1.25\times10^{-4}
\]

sugerindo que o sistema possui escalas temporais críticas
associadas ao regime dinâmico.

A observação não estabelece uma lei fundamental, mas indica
uma possível relação entre intensidade relacional e resolução
temporal necessária para preservar estabilidade.
\section{Auditoria 15 — Lei de Escala Temporal Crítica}

A curva de estabilidade temporal foi analisada através de um
ajuste em escala logarítmica.

Assumindo uma relação fenomenológica:

\[
\Delta t_{crit}=C\beta^{-\alpha}
\]

foi realizado um ajuste linear no espaço:

\[
\log(\Delta t_{crit})
=
\log(C)-\alpha\log(\beta)
\]

Os resultados obtidos foram:

\[
\alpha=0.41696
\]

\[
C=5.5049\times10^{-4}
\]

com coeficiente de determinação:

\[
R^2=0.9270
\]

O resultado indica uma dependência sistemática entre a
intensidade não linear $\beta$ e a escala temporal crítica.

A relação observada é compatível com uma lei de potência
aproximada:

\[
\Delta t_{crit}
\propto
\beta^{-0.417}
\]

Entretanto, este resultado permanece como uma caracterização
empírica do modelo S26-B, devendo ser submetido a testes com
maior faixa de parâmetros e diferentes topologias.
\section{Auditoria 16 — Transição de Regime na Escala Temporal Crítica}

A análise refinada da dependência entre o parâmetro não linear
$\beta$ e o passo temporal crítico revelou dois regimes distintos.

Para baixos valores de $\beta$, observou-se um platô de estabilidade:

\[
\Delta t_{crit}\approx constante
\]

indicando que a limitação temporal não depende fortemente da
não linearidade nesse intervalo.

Entretanto, para o regime de maior intensidade não linear,
o ajuste indicou:

\[
\Delta t_{crit}\propto\beta^{-0.495}
\]

comportamento próximo de:

\[
\Delta t_{crit}\propto\beta^{-1/2}
\]

Os resultados sugerem uma transição entre um regime dominado
pela estabilidade estrutural e outro dominado pela escala
dinâmica não linear.
\section{Auditoria 17 — Localização da Transição de Estabilidade}

A varredura refinada do parâmetro não linear $\beta$ revelou uma
mudança no limite temporal de estabilidade.

Para:

\[
\beta \leq 12
\]

observou-se:

\[
\Delta t_{crit}=2.5\times10^{-4}
\]

Enquanto para:

\[
\beta \geq 14
\]

o limite observado tornou-se:

\[
\Delta t_{crit}=1.25\times10^{-4}
\]

A região de transição encontra-se aproximadamente em:

\[
12 < \beta_c < 14
\]

indicando uma possível mudança de regime dinâmico.
\section{Auditoria 18 — Refinamento da Fronteira Crítica}

A investigação da região próxima ao ponto de transição
previamente identificado revelou que a mudança no limite temporal
não ocorre de forma abrupta.

A análise refinada:

\[
12\leq\beta\leq14
\]

mostrou manutenção do limite:

\[
\Delta t_{crit}=2.5\times10^{-4}
\]

até aproximadamente:

\[
\beta\approx13.5
\]

seguida por uma redução gradual:

\[
\Delta t_{crit}=2.0\times10^{-4}
\]

em:

\[
\beta=14
\]

O resultado indica um regime de crossover entre estabilidade
estrutural e limitação dinâmica, em vez de uma transição
descontínua.
\section{Auditoria 19 — Mapeamento do Crossover Temporal}

A região de transição entre o regime de estabilidade estrutural
e o regime dominado pela não linearidade foi investigada através
de uma varredura refinada.

Os resultados mostraram uma redução gradual do passo temporal
crítico:

\[
\Delta t_{crit}=2.25\times10^{-4}
\]

para:

\[
14\leq\beta\leq16
\]

e:

\[
\Delta t_{crit}=2.0\times10^{-4}
\]

para:

\[
17\leq\beta\leq20
\]

A ausência de uma descontinuidade indica um crossover dinâmico
em vez de uma transição abrupta.

O comportamento sugere a existência de uma região intermediária
entre o regime de baixa não linearidade e o regime assintótico de
alta intensidade.
\section{Auditoria 13 — Mapa de Estabilidade Temporal}

A estabilidade do integrador foi investigada em função da intensidade
não linear $\beta$.

Inicialmente observou-se um regime de estabilidade praticamente
independente de $\beta$, seguido por uma região de transição e,
posteriormente, por um regime assintótico no qual o passo temporal
crítico diminui sistematicamente com o aumento da não linearidade.

Os experimentos permitiram identificar três regiões distintas:

\begin{enumerate}

\item Regime de platô ($\beta \lesssim 13.5$);

\item Regime de transição ($14 \lesssim \beta \lesssim 25$);

\item Regime assintótico ($\beta \gtrsim 25$).

\end{enumerate}

No regime assintótico foi ajustada uma lei de potência da forma

\[
\Delta t_{\mathrm{crit}}
=
C
\beta^{-\alpha}.
\]

Foram obtidos os seguintes resultados experimentais:

\begin{table}[h]
\centering
\begin{tabular}{ccc}
\toprule
Intervalo & $\alpha$ & $R^2$\\
\midrule
$25 \leq \beta \leq 100$ & 0.483 & 0.932\\
$\beta \geq 30$ & 0.535 & 0.941\\
\bottomrule
\end{tabular}
\end{table}

Os dois ajustes apresentam expoentes próximos de $1/2$, indicando que,
na região fortemente não linear investigada, o passo temporal crítico é
compatível com uma dependência aproximadamente proporcional a

\[
\Delta t_{\mathrm{crit}}
\propto
\beta^{-1/2}.
\]

Este comportamento foi obtido exclusivamente a partir das simulações
numéricas do modelo e caracteriza o regime assintótico observado para a
implementação atual.

A etapa de validação numérica foi concluída com a caracterização do
domínio de estabilidade temporal do integrador, incluindo a identificação
de um regime de platô, uma região de transição e um regime assintótico
compatível com uma lei de potência aproximadamente proporcional a
$\Delta t_{\mathrm{crit}}\propto\beta^{-1/2}$.

Com esta infraestrutura validada, as próximas auditorias concentram-se
na caracterização espectral da dinâmica relacional, iniciando a Fase B2
do projeto.
% ------------------------------------------------
\section{S26-B34.1 — Robustez Temporal das Leis de Mistura Espectral}

A auditoria S26-B34.1 foi realizada para verificar se as leis
de escala obtidas anteriormente na S26-B33 representam uma
propriedade intrínseca do sistema ou um efeito associado à
discretização temporal.

Foram recalculados os expoentes de crescimento das taxas máximas
de mistura espectral para diferentes passos temporais:

\[
\Delta t =
2.50\times10^{-4},
1.25\times10^{-4},
6.25\times10^{-5},
3.125\times10^{-5}
\]

considerando as relações de escala:

\[
\max\left(\frac{dS}{dt}\right)
\propto \beta^{\alpha_S}
\]

\[
\max\left(\frac{dW}{dt}\right)
\propto \beta^{\alpha_W}
\]

\[
\max\left(\frac{dA}{dt}\right)
\propto \beta^{\alpha_A}
\]

Os resultados obtidos foram:

\begin{table}[h]
\centering
\begin{tabular}{cccc}
\toprule
$\Delta t$ &
$\alpha_S$ &
$\alpha_W$ &
$\alpha_A$
\\
\midrule
$2.50e-4$ &
2.2227 &
1.7034 &
1.8755
\\

$1.25e-4$ &
1.1847 &
1.0625 &
1.1275
\\

$6.25e-5$ &
1.0450 &
1.0278 &
1.0288
\\

$3.13e-5$ &
1.0115 &
1.0263 &
1.0070
\\
\bottomrule
\end{tabular}
\end{table}


Observa-se uma convergência sistemática dos expoentes conforme
o refinamento temporal.

Os valores inicialmente superlineares encontrados em passos
temporais maiores não permanecem no limite de resolução temporal
adequada, indicando que representam efeitos de sub-resolução
numérica.


No regime temporalmente convergido, os expoentes aproximam-se de:

\[
\alpha_S \approx 1
\]

\[
\alpha_W \approx 1
\]

\[
\alpha_A \approx 1
\]


Portanto, a caracterização corrigida da dinâmica espectral indica
crescimento aproximadamente linear das taxas de mistura com o
parâmetro não linear:

\[
\max\left(\frac{dS}{dt}\right)
\sim \beta
\]

\[
\max\left(\frac{dW}{dt}\right)
\sim \beta
\]

\[
\max\left(\frac{dA}{dt}\right)
\sim \beta
\]


Este resultado substitui a interpretação preliminar da S26-B33,
demonstrando que auditorias de convergência são essenciais para
distinguir leis dinâmicas reais de efeitos introduzidos pela
discretização numérica.


\subsection{Conclusão parcial}

A auditoria S26-B34.1 confirma que a intensidade da mistura
espectral possui comportamento aproximadamente linear com a
não linearidade no regime resolvido numericamente.

As próximas etapas da S26-B34 investigarão a robustez espacial
e a dependência com a condição inicial, variando respectivamente
o tamanho da rede e a largura inicial do pacote.

% ------------------------------------------------
\section{S26-B34.2 — Robustez Espacial das Leis de Escala}

Após a validação temporal realizada na S26-B34.1, foi investigada
a dependência das leis de escala em relação à resolução espacial
da rede relacional.

Foram utilizados diferentes tamanhos de rede:

\[
n = 128,\ 256,\ 512
\]

mantendo o regime temporal convergido:

\[
\Delta t = 6.25\times10^{-5}
\]

Os expoentes obtidos foram:

\[
\begin{array}{c|ccc}
n &
\alpha_S &
\alpha_W &
\alpha_A
\\
\hline

128 &
1.0456 &
1.0285 &
1.0288
\\

256 &
1.0450 &
1.0278 &
1.0288
\\

512 &
1.0448 &
1.0095 &
1.0288
\\

\end{array}
\]

Observa-se estabilidade dos expoentes com o aumento da resolução
espacial.

A entropia espectral e a amplitude apresentam variação inferior ao
limite esperado para discretização numérica, enquanto a largura
modal converge para valores próximos da unidade.

Portanto, a dependência aproximadamente linear:

\[
\alpha_S\approx\alpha_W\approx\alpha_A\approx1
\]

permanece após refinamento espacial.

Este resultado confirma que a lei de escala observada não é um
artefato específico do tamanho da rede utilizada.
% ------------------------------------------------
\section{S26-B34 — Auditoria de Robustez das Leis de Mistura Espectral}

A etapa S26-B34 teve como objetivo verificar se as leis de escala
observadas na caracterização espectral dependem de parâmetros
numéricos ou de escolhas específicas da simulação.

Foram realizadas três auditorias independentes:

\begin{itemize}

\item robustez temporal;

\item robustez espacial;

\item robustez da condição inicial.

\end{itemize}


\subsection{S26-B34.1 — Robustez Temporal}

Inicialmente, a S26-B33 apresentou expoentes superiores à unidade
para as taxas de mistura espectral.

Para verificar se esse comportamento era estrutural ou um efeito
de discretização, foram realizados refinamentos sucessivos do
passo temporal.

Os expoentes obtidos foram:

\[
\begin{array}{c|ccc}
\Delta t &
\alpha_S &
\alpha_W &
\alpha_A
\\
\hline

2.50e-4 &
2.2227 &
1.7034 &
1.8755
\\

1.25e-4 &
1.1847 &
1.0625 &
1.1275
\\

6.25e-5 &
1.0450 &
1.0278 &
1.0288
\\

3.13e-5 &
1.0115 &
1.0263 &
1.0070
\\

\end{array}
\]

Observou-se convergência dos expoentes para valores próximos da
unidade quando o regime temporal foi resolvido adequadamente.

Portanto, os valores superlineares iniciais foram classificados
como efeitos de sub-resolução temporal.


\subsection{S26-B34.2 — Robustez Espacial}

Após a convergência temporal, foi investigada a dependência das
leis de escala com o tamanho da rede.

Foram utilizados:

\[
n=128,\ 256,\ 512
\]

mantendo:

\[
\Delta t=6.25\times10^{-5}
\]

Os resultados foram:

\[
\begin{array}{c|ccc}
n &
\alpha_S &
\alpha_W &
\alpha_A
\\
\hline

128 &
1.0456 &
1.0285 &
1.0288
\\

256 &
1.0450 &
1.0278 &
1.0288
\\

512 &
1.0448 &
1.0095 &
1.0288
\\

\end{array}
\]

A variação espacial apresentou impacto desprezível nos expoentes,
indicando independência da discretização espacial no regime testado.


\subsection{S26-B34.3 — Robustez da Condição Inicial}

Finalmente, foi avaliada a influência da largura inicial do pacote
relacional.

Foram utilizados:

\[
\sigma=10,\ 20,\ 40
\]

com os demais parâmetros mantidos constantes.

Os resultados foram:

\[
\begin{array}{c|ccc}
\sigma &
\alpha_S &
\alpha_W &
\alpha_A
\\
\hline

10 &
1.0401 &
1.0274 &
1.0288
\\

20 &
1.0450 &
1.0278 &
1.0288
\\

40 &
1.0473 &
1.0280 &
1.0288
\\

\end{array}
\]


A estabilidade dos expoentes frente à variação da condição inicial
indica que a lei observada não depende de uma configuração inicial
específica.


\subsection{Conclusão da Auditoria S26-B34}

As três auditorias realizadas demonstram que a dependência observada
das taxas de mistura espectral permanece estável frente às principais
fontes de variação numérica investigadas.

No regime resolvido:

\[
\alpha_S\approx\alpha_W\approx\alpha_A\approx1
\]

resultando em uma dependência aproximadamente linear:

\[
\frac{dM}{dt}\sim\beta
\]

onde $M$ representa as métricas espectrais monitoradas.

A S26-B34 transforma a observação preliminar da S26-B33 em uma
caracterização robusta, eliminando interpretações associadas à
discretização temporal, espacial ou à condição inicial.

Com esta etapa concluída, a Fase B2 — Caracterização Espectral —
é considerada encerrada.

A próxima fase investigará a existência de estados persistentes,
auto-localização e possíveis estados estacionários do operador
relacional.

\section{S26-B35 — Observatório de Persistência Dinâmica}

\subsection{Motivação}

Após a conclusão das etapas de validação numérica do motor dinâmico
(S26-V) e da caracterização espectral do regime não linear, inicia-se
a investigação sistemática de regimes persistentes no espaço de estados.

As etapas anteriores estabeleceram a confiabilidade do integrador,
a conservação das quantidades globais dentro da precisão numérica
esperada e a capacidade do observatório espectral de acompanhar a
redistribuição modal da dinâmica.

A presente etapa não busca identificar fenômenos específicos, mas
classificar matematicamente a evolução temporal do campo relacional
$\Gamma$ através de invariantes e funcionais quantitativos.

O objetivo é responder:

\[
\text{A trajetória dinâmica permanece transitória, recorrente ou persistente?}
\]

sem introduzir hipóteses externas ao formalismo.

% ----------------------------------------------------------

\subsection{Formulação Matemática}

O sistema investigado permanece definido pelo operador dinâmico:

\[
\ddot{\Gamma}
=
-\alpha L\Gamma
+
\beta F(\Gamma)
\]

onde:

\begin{itemize}
    \item $\Gamma \in \mathbb{R}^{n}$ representa o vetor de estado;
    \item $L$ representa o operador relacional discreto;
    \item $F(\Gamma)$ representa o termo não linear;
    \item $\beta$ controla a intensidade do acoplamento não linear.
\end{itemize}

A evolução numérica produz uma sequência discreta de estados:

\[
\Gamma(t_0),
\Gamma(t_1),
...
\Gamma(t_m)
\]

sobre a qual são definidos observáveis temporais e espectrais.

% ----------------------------------------------------------

\subsection{Observáveis de Persistência}

\subsubsection{Resíduo Local Temporal}

A primeira métrica mede a variação instantânea normalizada do campo:

\[
R_{loc}(t)
=
\frac{
\|\Gamma(t+\Delta t)-\Gamma(t)\|_2
}{
\Delta t\|\Gamma(t)\|_2+\epsilon
}
\]

Valores reduzidos indicam pequena evolução local do vetor de estado,
enquanto valores elevados indicam regimes de rápida transformação.

O observável não classifica o comportamento como estável ou instável
isoladamente, sendo utilizado em conjunto com os demais funcionais.

% ----------------------------------------------------------

\subsubsection{Divergência Espectral}

A distribuição modal normalizada é definida como:

\[
p_k(t)
=
\frac{
|\hat{\Gamma}_k(t)|^2
}{
\sum_j |\hat{\Gamma}_j(t)|^2
}
\]

A variação espectral entre passos consecutivos é medida pela
divergência Jensen-Shannon:

\[
D_{spec}
=
JS
(
p(t),
p(t+\Delta t)
)
\]

Essa métrica quantifica a reorganização da energia no espaço espectral.

% ----------------------------------------------------------

\subsubsection{Evolução da Forma Espectral}

A distribuição espacial da energia modal é acompanhada pelo
Participation Ratio:

\[
PR(t)
=
\frac{
1
}{
\sum_k p_k(t)^2
}
\]

A variação temporal é definida como:

\[
H_{shape}
=
\frac{
PR(t+\Delta t)-PR(t)
}{
\Delta t
}
\]

Este observável mede expansão ou contração efetiva da ocupação modal.

% ----------------------------------------------------------

\subsubsection{Coerência Temporal}

A memória do estado inicial é medida pela correlação:

\[
C_{auto}(t)
=
\frac{
\Gamma(0)\cdot\Gamma(t)
}{
\|\Gamma(0)\|_2
\|\Gamma(t)\|_2
}
\]

Valores próximos de unidade indicam preservação da direção inicial
do vetor de estado, enquanto oscilações ou redução indicam perda de
correlação temporal.

% ----------------------------------------------------------

\subsubsection{Entropia Espectral}

A entropia normalizada é definida por:

\[
S(t)
=
-\frac{
\sum_k p_k(t)\log(p_k(t))
}{
\log(n)
}
\]

Este funcional quantifica a distribuição relativa da energia entre os
modos disponíveis.

A entropia é interpretada exclusivamente como medida de dispersão
espectral.

% ----------------------------------------------------------

\subsubsection{Distância de Energia Modal}

Como complemento à divergência espectral, define-se:

\[
D_E(t)
=
\left\|
E_k(t+\Delta t)
-
E_k(t)
\right\|_2
\]

onde:

\[
E_k(t)
=
\frac{
|\hat{\Gamma}_k(t)|^2
}{
\sum_j |\hat{\Gamma}_j(t)|^2
}
\]

Este observável fornece uma medida direta da alteração da distribuição
energética modal.

% ----------------------------------------------------------

\subsection{Classificação de Regimes Dinâmicos}

Os observáveis anteriores permitem uma classificação puramente
algorítmica:

\begin{itemize}

\item \textbf{Regime Transitório:}
as métricas permanecem com variação significativa no tempo.

\item \textbf{Regime Persistente:}
as métricas coletivas apresentam estabilização dentro de uma tolerância
definida.

\item \textbf{Regime Oscilatório:}
as métricas permanecem limitadas, porém apresentam recorrência temporal.

\item \textbf{Regime Instável:}
as métricas apresentam crescimento não controlado ou divergência
numérica.

\end{itemize}

Nenhuma dessas classificações implica interpretação fenomenológica.

% ----------------------------------------------------------

\subsection{Critério Metodológico}

A etapa S26-B35 estabelece um observatório de persistência dinâmica.

O objetivo é determinar propriedades internas da trajetória no espaço
de estados:

\[
\Gamma(t)
\rightarrow
\{\mathcal{M}_1(t),
\mathcal{M}_2(t),
...,
\mathcal{M}_n(t)\}
\]

onde cada $\mathcal{M}$ representa um funcional mensurável do sistema.

A análise permanece restrita ao formalismo discreto:

\begin{itemize}
    \item vetor de estado em $\mathbb{R}^{n}$;
    \item operador relacional discreto;
    \item espectro do operador;
    \item invariantes temporais e energéticos.
\end{itemize}

A interpretação dos resultados será realizada somente após a auditoria
dos dados produzidos pelo observatório.

% ==========================================================
\subsection{Extensão do Observatório: Centro Modal}

Como complemento ao conjunto de observáveis de persistência foi
introduzido o resíduo de deslocamento do centro modal.

Seja

\[
C_m(t)
\]

o centro modal calculado pelo observatório espectral.

Define-se

\[
R_{macro}
=
\frac{
|C_m(t+\Delta t)-C_m(t)|
}{
n\,\Delta t
}
\]

onde \(n\) representa o número de graus de liberdade da malha.

Este observável mede exclusivamente a migração global do centro da
distribuição modal, permanecendo independente das alterações internas da
largura espectral.

Dessa forma torna-se possível distinguir entre:

\begin{itemize}

\item redistribuição interna da energia modal;

\item translação global do espectro.

\end{itemize}

Os dois processos são matematicamente independentes.

% ----------------------------------------------------------

\subsection{Auditoria Numérica do Observatório}

O conjunto de observáveis foi aplicado inicialmente ao regime
\(\beta=1\).

Obtiveram-se os seguintes resultados médios:

\[
\langle R_{loc}\rangle
=
7.56\times10^{-2}
\]

\[
\langle D_{spec}\rangle
=
2.71\times10^{-11}
\]

\[
\langle H_{shape}\rangle
=
3.51\times10^{-1}
\]

\[
\langle C_{auto}\rangle
=
0.99999999995
\]

\[
\langle R_{macro}\rangle
=
8.83\times10^{-4}
\]

A entropia espectral permaneceu praticamente constante durante toda a
janela observada, apresentando desvio padrão da ordem de

\[
10^{-6}.
\]

% ----------------------------------------------------------

\subsection{Consistência Interna}

Os observáveis apresentaram comportamento mutuamente consistente.

Observou-se simultaneamente:

\begin{itemize}

\item evolução temporal não nula do vetor de estado
(\(R_{loc}>0\));

\item divergência espectral praticamente nula;

\item deslocamento praticamente nulo do centro modal;

\item manutenção quase perfeita da autocorrelação temporal;

\item estabilidade da entropia espectral.

\end{itemize}

Esses resultados indicam que, dentro da janela temporal investigada,
a trajetória permanece confinada a uma região extremamente restrita do
espaço de estados, preservando praticamente inalterada sua estrutura
modal global.

Tal comportamento não constitui demonstração da existência de estados
estacionários, mas caracteriza um candidato para investigação
sistemática na etapa S26-B36.

% ----------------------------------------------------------

\subsection{Conclusão da Etapa S26-B35}

A etapa S26-B35 estabelece um observatório permanente de persistência
dinâmica baseado em observáveis independentes definidos sobre o espaço
físico e o espaço espectral.

O módulo passa a integrar permanentemente a biblioteca do projeto,
fornecendo os funcionais utilizados pelas etapas posteriores de
classificação dinâmica.

A etapa é considerada encerrada após a validação numérica dos
observáveis e sua incorporação definitiva ao ambiente experimental.

%=========================================================
\subsection{S26-B36 — Auditoria do Classificador de Regimes}

A etapa S26-B36 introduziu a primeira camada de classificação
automática dos regimes dinâmicos produzidos pelo Observatório de
Persistência desenvolvido na etapa S26-B35.

O objetivo desta fase não consiste em descobrir novos fenômenos
físicos, mas em verificar se os observáveis coletivos são suficientes
para distinguir diferentes comportamentos assintóticos da dinâmica.

A classificação é realizada exclusivamente a partir dos observáveis
internos calculados durante a simulação, sem utilização de parâmetros
externos ou limiares fenomenológicos.

%---------------------------------------------------------

\subsubsection{Critério de Classificação}

A decisão passa a ser baseada sobre uma janela terminal da evolução
numérica.

São utilizados:

\[
\langle R_{loc}\rangle,
\]

\[
\langle D_{spec}\rangle,
\]

\[
\langle |\Delta PR| \rangle,
\]

\[
Var(R_{loc}),
\]

\[
Var(D_{spec}),
\]

bem como as inclinações temporais das métricas coletivas.

Optou-se por utilizar diretamente

\[
\Delta PR
=
PR(t+\Delta t)-PR(t),
\]

em substituição ao funcional

\[
H_{shape}
=
\frac{\Delta PR}{\Delta t},
\]

na etapa de classificação.

Essa alteração elimina a dependência artificial introduzida pelo passo
temporal, preservando apenas a variação física efetiva do suporte
espectral.

%---------------------------------------------------------

\subsubsection{Persistence Score}

O índice agregado de persistência permanece definido por

\[
P(t)
=
|C_{auto}(t)|
\exp
\left[
-
\left(
\Delta t\,R_{loc}
+
D_{spec}
+
|\Delta PR|
\right)
\right].
\]

Esta forma apresenta as propriedades:

\begin{itemize}

\item
\(0\le P\le1\);

\item
continuidade;

\item
invariância de escala;

\item
ausência de pesos arbitrários.

\end{itemize}

Dessa forma o escore resume a estabilidade coletiva da trajetória em um
único funcional escalar.

%---------------------------------------------------------

\subsubsection{Primeira Auditoria Numérica}

Aplicando o classificador ao conjunto produzido pelo observatório
S26-B35 obteve-se:

\[
P
=
0.99995731040597,
\]

com variância

\[
Var(P)
=
4.03\times10^{-10}.
\]

As médias dos principais observáveis foram

\[
\langle R_{loc}\rangle
=
7.56\times10^{-2},
\]

\[
\langle D_{spec}\rangle
=
2.71\times10^{-11},
\]

\[
\langle |\Delta PR|\rangle
=
3.51\times10^{-5}.
\]

O coeficiente de autocorrelação permaneceu praticamente unitário ao
longo de toda a janela terminal.

%---------------------------------------------------------

\subsubsection{Interpretação}

A auditoria evidencia que o sistema permanece extremamente próximo de um
estado estacionário coletivo.

Entretanto,

\[
R_{loc}\neq0,
\]

mostrando que pequenas variações microscópicas continuam presentes.

Consequentemente, o termo
\emph{persistente}
não significa ausência completa de evolução, mas sim estabilidade das
coordenadas coletivas do sistema.

A distinção entre estabilidade macroscópica e flutuação microscópica
passa a constituir um dos principais resultados metodológicos da série
S26.

%---------------------------------------------------------

\subsubsection{Resultado}

O classificador identificou automaticamente o regime

\[
\boxed{\text{PERSISTENTE}}
\]

para a dinâmica analisada.

Este resultado valida simultaneamente:

\begin{itemize}

\item
o Observatório de Persistência (S26-B35);

\item
o Persistence Score;

\item
a árvore de decisão implementada na etapa S26-B36.

\end{itemize}

A próxima etapa investigará a robustez desta classificação sob
variações de resolução numérica, tamanho da malha e intensidade do
parâmetro não linear, estabelecendo o mapa completo de regimes
dinâmicos do sistema GER.

%%%%%%%%%%%%%%%%%%%%%%%%%%%%%%%%%%%%%%%%%%%%%%%%%%%%%%%%%%%%%%
\section{S26-B36 — Reconstrução Científica do Stationary Scan}

A reconstrução do módulo S26-B36 foi conduzida como uma revisão conceitual completa da
etapa originalmente denominada ``Classificador de Regimes''.

As auditorias numéricas, a análise dos observáveis produzidos pelo S26-B35 e a formulação
matemática desenvolvida durante esta reconstrução indicaram que o objeto de estudo do
módulo não deve ser um observável individual nem o vetor de observáveis instantâneo.

O objeto matemático do S26-B36 passa a ser a trajetória produzida pelo observatório.

%%%%%%%%%%%%%%%%%%%%%%%%%%%%%%%%%%%%%%%%%%%%%%%%%%%%%%%%%%%%%%
\subsection{Objeto Matemático}

Seja

\[
x(t)\in\mathcal O
\]

o vetor de observáveis produzido pelo S26-B35.

Define-se a trajetória observacional

\[
\gamma=\{x(t):t\in[t_0,t_f]\}\subset\mathcal O.
\]

O módulo S26-B36 passa a estudar propriedades geométricas globais da trajetória
\(\gamma\), e não propriedades de um estado isolado.

Consequentemente, um regime dinâmico deixa de ser interpretado como um atributo de um
observável individual e passa a representar uma propriedade geométrica da trajetória no
espaço de observação.

%%%%%%%%%%%%%%%%%%%%%%%%%%%%%%%%%%%%%%%%%%%%%%%%%%%%%%%%%%%%%%
\section{GER-S26-B36-OGF}

\subsection{Operadores Geométricos Fundamentais}

O documento GER-S26-B36-OGF estabelece a interface matemática oficial entre o
Observatório de Persistência (S26-B35) e o Stationary Scan.

Todo operador geométrico fundamental do S26-B36 é definido como um funcional sobre a
trajetória observacional

\[
\mathcal F[\gamma].
\]

Os operadores fundamentais são:

\begin{itemize}
\item Confinamento;
\item Convergência;
\item Recorrência;
\item Deriva.
\end{itemize}

Esses operadores constituem a linguagem geométrica oficial utilizada pelo módulo.

%%%%%%%%%%%%%%%%%%%%%%%%%%%%%%%%%%%%%%%%%%%%%%%%%%%%%%%%%%%%%%
\subsection{Arquitetura Conceitual}

A arquitetura científica do módulo passa a ser

\[
\text{Trajetória}
\longrightarrow
\text{Operadores Geométricos}
\longrightarrow
\text{Diagnóstico Geométrico}
\longrightarrow
\text{Classificação}.
\]

O diagnóstico produzido pelo Stationary Scan constitui uma inferência sobre a geometria da
trajetória e não uma classificação direta dos operadores individualmente.

Assim, nenhum operador isolado caracteriza um regime.

Os regimes emergem da interpretação conjunta dos operadores geométricos fundamentais.

%%%%%%%%%%%%%%%%%%%%%%%%%%%%%%%%%%%%%%%%%%%%%%%%%%%%%%%%%%%%%%
\subsection{Consequência Metodológica}

O State Vector permanece como parte integrante do observatório.

Entretanto, sua função passa a ser exclusivamente descritiva.

As estatísticas produzidas pelo S26-B35 resumem propriedades distintas da trajetória, mas
não constituem, individualmente, critérios classificatórios.

%%%%%%%%%%%%%%%%%%%%%%%%%%%%%%%%%%%%%%%%%%%%%%%%%%%%%%%%%%%%%%
\section{Marco Científico --- Reconstrução do S26-B36}

A reconstrução científica do S26-B36 foi considerada concluída com o congelamento do
documento GER-S26-B36-OGF (Operadores Geométricos Fundamentais).

A partir desta decisão, o objeto matemático oficial do módulo passa a ser a trajetória
observacional produzida pelo S26-B35, e não mais observáveis individuais ou o State
Vector instantâneo.

Os operadores fundamentais (Confinamento, Convergência, Recorrência e Deriva)
constituem a interface matemática oficial entre o Observatório e o Stationary Scan.

Este documento estabelece a especificação matemática oficial do módulo, vigente até
eventual revisão motivada por evidências experimentais reproduzíveis.

%%%%%%%%%%%%%%%%%%%%%%%%%%%%%%%%%%%%%%%%%%%%%%%%%%%%%%%%%%%%%%
\subsection{Princípio de Evolução Científica}

A evolução futura do S26-B36 deverá ocorrer prioritariamente pela ampliação do domínio
experimental e pela validação dos Operadores Geométricos Fundamentais.

Alterações na teoria somente deverão ser consideradas quando evidências experimentais

reproduzíveis demonstrarem incompatibilidade entre os operadores definidos no
GER-S26-B36-OGF e o comportamento observado do sistema.

Na ausência dessas evidências, modificações na implementação, nos algoritmos de
aproximação ou nos parâmetros de calibração não constituem alterações da teoria.

%%%%%%%%%%%%%%%%%%%%%%%%%%%%%%%%%%%%%%%%%%%%%%%%%%%%%%%%%%%%%%
\subsection{Considerações Finais}

A reconstrução do S26-B36 estabeleceu uma separação explícita entre observação,
diagnóstico geométrico e classificação.

A arquitetura científica do projeto passa a ser

\[
\text{Motor}
\longrightarrow
\text{Observatório}
\longrightarrow
\text{Trajetória}
\longrightarrow
\text{Operadores Geométricos}
\longrightarrow
\text{Diagnóstico Geométrico}
\longrightarrow
\text{Classificação}.
\]

Essa organização torna independentes a teoria, sua implementação computacional e sua
calibração experimental, estabelecendo uma base metodológica para a evolução futura do
GER.

A partir deste marco, o desenvolvimento do S26-B36 deixa de consistir na descoberta de
uma teoria e passa a concentrar-se na implementação, validação experimental e
aperfeiçoamento computacional da teoria geométrica estabelecida pelo documento
GER-S26-B36-OGF.

\section{S26-B36 --- Implementação e Validação do Stationary Scan}

Após o congelamento da especificação matemática estabelecida pelo documento
GER-S26-B36-OGF, o desenvolvimento do módulo passou a concentrar-se
exclusivamente na implementação computacional da arquitetura definida.

O operador $\Psi$ foi implementado como o núcleo do \textit{Stationary Scan},
recebendo exclusivamente a Assinatura Geométrica

\[
\mathcal{G}
=
(\mathrm{diameter},
\mathrm{convergence},
\mathrm{recurrence},
\mathrm{drift})
\]

e produzindo um Certificado Estrutural de acordo com o axioma SS-1.

Durante esta etapa foram implementados os três operadores internos da
arquitetura:

\begin{itemize}
\item Eligibility Operator;
\item Structural Operator;
\item Consistency Operator.
\end{itemize}

O Eligibility Operator verifica a integridade da Assinatura Geométrica e a
elegibilidade das relações fundamentais.

O Structural Operator foi implementado como parte permanente da arquitetura.
Entretanto, durante a implementação e auditoria verificou-se que a teoria
vigente não fornece, até o presente momento, nenhuma dedução estrutural
formalmente demonstrada utilizando exclusivamente os Operadores
Geométricos Fundamentais.

Consequentemente, o operador permanece implementado, porém sem Rules
Estruturais ativas.

Esta decisão preserva rigorosamente o axioma SS-1, impedindo que o módulo
introduza conhecimento não demonstrado pela teoria.

O Consistency Operator verifica a consistência interna do Certificado
Estrutural, garantindo compatibilidade entre a implementação computacional
e o estado atual da teoria.

\subsection{Auditoria Oficial}

Foi desenvolvido um módulo independente de auditoria destinado a verificar
a conformidade da implementação com a especificação oficial do
Stationary Scan.

A auditoria valida:

\begin{itemize}
\item estrutura do Certificado Estrutural;
\item construção das seis relações fundamentais;
\item funcionamento do Eligibility Operator;
\item consistência do Structural Operator;
\item funcionamento do Consistency Operator;
\item reprodutibilidade determinística do Certificado.
\end{itemize}

Todos os testes foram aprovados sem inconsistências.

A implementação apresentou comportamento determinístico e reproduzível,
produzindo certificados idênticos para assinaturas geométricas idênticas.

\subsection{Conclusão da Etapa}

A implementação computacional do módulo S26-B36 é considerada concluída.

O operador $\Psi$ permanece arquiteturalmente congelado, estando preparado
para receber futuras Rules Estruturais formalmente demonstradas sem que
seja necessária qualquer alteração em sua arquitetura.

Na ausência dessas Rules, o Structural Operator permanece corretamente
implementado e cientificamente consistente com a teoria vigente.

Com esta etapa encerrada, considera-se concluída a reconstrução científica,
a implementação e a validação do S26-B36.

A próxima etapa do projeto inicia o módulo S26-B37, responsável por utilizar
o Certificado Estrutural produzido pelo Stationary Scan como entrada para a
camada subsequente da arquitetura da Geometria Espectral Relacional.
\section{Início do S26-B37 --- Geometria do Espaço das Assinaturas}

Com a conclusão e validação oficial do módulo S26-B36, iniciou-se a investigação
científica do módulo S26-B37.

Diferentemente da etapa anterior, o objetivo inicial não consistiu em propor
novos operadores dedutivos, mas em investigar propriedades fundamentais do
próprio espaço das Assinaturas Geométricas.

A primeira hipótese considerada foi a possibilidade de definir relações de
equivalência estrutural entre assinaturas. Entretanto, verificou-se que essa
questão depende previamente da existência de uma geometria bem definida sobre o
espaço

\[
\mathcal{G}
=
(\mathrm{diameter},
\mathrm{convergence},
\mathrm{recurrence},
\mathrm{drift}).
\]

Consequentemente, o primeiro objetivo do S26-B37 passou a ser a investigação
experimental da geometria induzida pelas diferentes métricas possíveis sobre
esse espaço.

\subsection{Experimento S26-B37-0}

Foi implementado o primeiro experimento destinado à validação da infraestrutura
experimental do módulo.

O experimento gerou um conjunto de assinaturas geométricas sintéticas e
verificou:

\begin{itemize}
\item geração das assinaturas;
\item representação canônica;
\item detecção de colisões;
\item cálculo de distâncias;
\item geração automática de relatórios.
\end{itemize}

Para um conjunto de 1000 assinaturas observou-se:

\begin{itemize}
\item 1000 assinaturas distintas;
\item nenhuma colisão;
\item infraestrutura computacional validada.
\end{itemize}

Concluiu-se que o experimento validou a infraestrutura necessária para os
experimentos subsequentes, não sendo ainda suficiente para estabelecer
propriedades geométricas da teoria.

\subsection{Experimento S26-B37-1}

O segundo experimento investigou a influência da métrica adotada sobre a
estrutura local do espaço das Assinaturas Geométricas.

Utilizando exatamente o mesmo conjunto de assinaturas, foram comparadas duas
métricas:

\begin{itemize}
\item distância euclidiana;
\item distância euclidiana normalizada.
\end{itemize}

Para cada métrica foi determinado o par de assinaturas mais próximo.

O resultado experimental mostrou que os pares obtidos pelas duas métricas foram
distintos.

Em outras palavras, a estrutura de vizinhança do espaço das Assinaturas depende
da métrica utilizada.

Este resultado constitui a primeira evidência experimental do S26-B37.

\subsection{Consequência Científica}

O experimento demonstra que a geometria do espaço das Assinaturas Geométricas
não pode ser assumida implicitamente como euclidiana.

Em consequência, qualquer teoria futura envolvendo:

\begin{itemize}
\item proximidade estrutural;
\item estabilidade;
\item continuidade;
\item equivalência;
\item dedutibilidade estrutural;
\end{itemize}

deverá ser formulada sobre uma métrica explicitamente definida.

O resultado não identifica qual é a métrica canônica do espaço das
Assinaturas, mas demonstra experimentalmente que essa escolha influencia a
estrutura geométrica observada.

\subsection{Questão Aberta}

O primeiro problema científico formal do S26-B37 passa a ser:

\begin{quote}
\emph{Existe uma métrica estrutural canônica para o espaço das Assinaturas
Geométricas produzidas pelos Operadores Geométricos Fundamentais?}
\end{quote}

A resposta a essa questão torna-se condição necessária para o desenvolvimento
das futuras teorias de equivalência estrutural, estabilidade dedutiva e
dedutibilidade no âmbito da Geometria Espectral Relacional.
\subsection{Experimento S26-B37-3 --- Geometria Global das Assinaturas}

Após a constatação de que a estrutura de vizinhança depende da métrica
adotada, foi realizado um experimento destinado a caracterizar a geometria
global induzida pela distância euclidiana sobre o espaço das Assinaturas
Geométricas.

Diferentemente dos experimentos anteriores, que analisavam propriedades
locais, este experimento avaliou todos os pares de assinaturas produzidos
durante a simulação, acumulando as contribuições relativas de cada Operador
Geométrico Fundamental para a distância média observada.

Foram executadas sucessivamente três campanhas experimentais contendo
1000, 2500 e 5000 assinaturas.

Os resultados convergiram para a mesma distribuição estatística, indicando
estabilidade numérica do observável.

A contribuição média observada foi aproximadamente:

\[
\begin{array}{c|c}
\text{Operador} & \text{Contribuição Média} \\
\hline
\mathrm{diameter} & \approx 1\% \\
\mathrm{convergence} & \approx 99\% \\
\mathrm{recurrence} & \approx 0\% \\
\mathrm{drift} & \approx 0\%
\end{array}
\]

A convergência dos resultados para diferentes tamanhos de amostra indica
que esta distribuição não constitui um artefato estatístico da geração
aleatória utilizada no experimento.

\subsection{Consequência Experimental}

O experimento demonstra que a métrica euclidiana induz uma geometria global
fortemente anisotrópica sobre o espaço das Assinaturas Geométricas.

Nesta geometria, a contribuição média do operador
\texttt{convergence} domina praticamente toda a distância entre
assinaturas, enquanto os demais Operadores Geométricos Fundamentais
apresentam contribuição global desprezível.

Este resultado não constitui uma propriedade dos Operadores Geométricos
Fundamentais, mas da métrica utilizada para comparar assinaturas.

Consequentemente, a escolha da métrica deixa de ser um detalhe de
implementação e passa a constituir parte integrante da formulação
matemática do módulo S26-B37.

\subsection{Problema Científico}

Os resultados experimentais obtidos conduzem ao seguinte problema
fundamental:

\begin{quote}
\emph{Construir uma métrica estrutural para o espaço das Assinaturas
Geométricas cuja geometria resulte da teoria dos Operadores Geométricos
Fundamentais e não apenas da escala numérica de seus observáveis.}
\end{quote}

Este problema passa a orientar o desenvolvimento das etapas seguintes do
S26-B37.

\subsection{Experimento S26-B37-6 --- Avaliação da Métrica Euclidiana Normalizada}

Após a caracterização da geometria global induzida pela distância
euclidiana, foi investigada uma primeira métrica candidata destinada a
eliminar a anisotropia observada.

A métrica considerada consiste na distância euclidiana aplicada às
componentes da Assinatura Geométrica após normalização individual de cada
Operador Geométrico Fundamental pela respectiva faixa observada durante o
experimento.

Mantendo exatamente o mesmo protocolo experimental utilizado no
Experimento S26-B37-3, foram realizadas quatro campanhas contendo
1000, 2500, 5000 e 10000 assinaturas geométricas sintéticas.

Para cada campanha foi calculada a contribuição média relativa de cada
Operador Geométrico Fundamental para a geometria global induzida pela
nova métrica.

Os resultados obtidos foram:

\[
\begin{array}{c|cccc}
\text{Amostras}
&
\mathrm{diameter}
&
\mathrm{convergence}
&
\mathrm{recurrence}
&
\mathrm{drift}
\\
\hline
1000
&
25.39\%
&
24.87\%
&
24.06\%
&
25.68\%
\\
2500
&
25.08\%
&
24.99\%
&
25.15\%
&
24.78\%
\\
5000
&
24.93\%
&
25.08\%
&
24.70\%
&
25.29\%
\\
10000
&
25.21\%
&
24.84\%
&
25.25\%
&
24.70\%
\end{array}
\]

Observa-se que os resultados convergem para uma distribuição
estatisticamente estável próxima de uma contribuição uniforme entre os
quatro Operadores Geométricos Fundamentais.

\subsection{Consequência Experimental}

Comparando diretamente os Experimentos S26-B37-3 e S26-B37-6, verifica-se
que a forte anisotropia observada sob a distância euclidiana desaparece
após a normalização das componentes da Assinatura Geométrica.

Os resultados obtidos indicam que a dominância anteriormente atribuída ao
operador \texttt{convergence} decorre predominantemente da diferença de
escala numérica entre os observáveis, e não necessariamente de uma
propriedade geométrica intrínseca da Assinatura.

Este resultado constitui uma evidência experimental robusta, reproduzida
para diferentes tamanhos de amostra, de que a escolha da métrica altera
substancialmente a geometria induzida sobre o espaço das Assinaturas.

\subsection{Estado Atual da Hipótese}

Os experimentos realizados permitem estabelecer que a métrica euclidiana
normalizada constitui uma candidata significativamente superior à métrica
euclidiana convencional quanto ao equilíbrio geométrico das contribuições
dos Operadores Geométricos Fundamentais.

Entretanto, os resultados atuais não permitem concluir que esta seja a
métrica estrutural definitiva do framework GER.

Permanecem em aberto, entre outras, as seguintes questões:

\begin{itemize}
\item comportamento sob distribuições não uniformes das Assinaturas;
\item comportamento para Assinaturas produzidas pelo motor físico do GER;
\item existência de métricas estruturalmente superiores à normalização
euclidiana.
\end{itemize}

Em consequência, a métrica euclidiana normalizada passa a constituir a
primeira métrica candidata oficialmente aceita para investigação nas
etapas subsequentes do módulo S26-B37, permanecendo sujeita às futuras
auditorias experimentais do framework.



\newpage
\input{journal}

\end{document}
